(1) 当回路仅部分进入磁场时，磁通量的纸内大小随重叠长度 $x$ 增大，$\Phi=Bwx$。由于 $dx/dt=v$，Faraday定律给电动势大小 $\mathcal E=|d\Phi/dt|=Bwv$。该式仅在一条竖边位于磁场内、另一边尚在外时成立。
(2) 通过欧姆定律得 $I=\mathcal E/R=Bwv/R$。进入时纸内磁通逐渐增大，感应电流应生成纸外磁场抵抗这一变化；依右手定则，从纸面正面看电流沿逆时针方向流动，符合Lenz定律。 电流方向也可用叉乘核验：纸外感应磁通对应逆时针环流，产生的磁作用力沿运动反方向，因而与电磁制动的能量收支相容。
(3) 有磁场作用的前侧竖边产生水平制动力 $F_B=IBw=B^2w^2v/R$，方向与运动速度相反。若匀速穿入，外力必须向右提供同样大小，机械输入功率为 $P_{\rm mech}=F_Bv=B^2w^2v^2/R$。
(4) 回路焦耳热率 $P_J=I^2R=(Bwv/R)^2R=B^2w^2v^2/R=P_{\rm mech}$，证明外机械功率全部转化成电阻热。回路全部置于匀强磁场后总磁通不再变化，各边动生电动势相抵消，理想模型下感应电流归零。
